% Part: many-valued-logic
% Chapter: syntax-and-semantics
% Section: valuations-sat

\documentclass[../../../include/open-logic-section]{subfiles}

\begin{document}

\olfileid[hi]{mvl}{syn}{val}

\olsection{\usetoken{P}{valuation} और संतुष्टि}

\begin{defn}[!!^{valuation}s] 
$V$ सत्य-मानों का कोई समुच्चय हो। $\Lang{L}$ से~$V$ में कोई
\emph{!!{valuation}} ऐसा फलन~$\pAssign{v}$ है जो भाषा के
!!{propositional variable} को~$V$ का !!a{element} नियत करता है, अर्थात्
$\pAssign{v} \colon \PVar \to V$।
\end{defn}

\begin{defn}\ollabel{defn:pValue}
  किसी बहुमूल्य तर्कशास्त्र~$\Log L$ के सत्य-मानों के समुच्चय~$V$ में
  !!a{valuation}~$\pAssign{v}$ दिए होने पर मूल्यांकन-फलन $\pValue{v}
  \colon \Frm[L] \to V$ को आगमनतः इस प्रकार परिभाषित करें:
  \begin{enumerate}
    \item $\pValue{v}(\Obj p_n) = \pAssign{v}(\Obj p_n)$; 
    \item यदि $\star$ कोई $0$-स्थानी संयोजक है, तो $\pValue{v}(\star)  = \tf{\star}[\Log L]$;
    \item यदि $\star$ कोई $n$-स्थानी संयोजक है, तो
    \[
      \pValue{v}(\star(!A_1, \dots, !A_n)) = \tf{\star}[\Log L]
      (\pValue{v}(!A_1), \dots, \pValue{v}(!A_n)).
    \]
  \end{enumerate}
\end{defn}

\begin{defn}[संतुष्टि]
\ollabel{defn:satisfaction} !!{formula}~$!A$ किसी
  !!a{valuation}~$\pAssign{v}$ द्वारा \emph{संतुष्ट} है,
  $\pSat{v}{!A}[\Log L]$, तभी और केवल तभी जब
  $\pValue{v}(!A)[\Log L] \in V^+$, जहाँ $V^+$, ~$\Log L$ के निर्दिष्ट
  सत्य-मानों का समुच्चय है।
  
  ``$\pSat{v}{!A}[\Log L]$ नहीं'' के अर्थ में हम
  $\pSat/{v}{!A}[\Log L]$ लिखते हैं। यदि $\Gamma$ !!{formula} का कोई
  समुच्चय है, तो $\pSat{v}{\Gamma}[\Log L]$ तभी और केवल तभी जब प्रत्येक
  ~$!A \in \Gamma$ के लिए $\pSat{v}{!A}[\Log L]$।
\end{defn}

\end{document}
